\documentclass[book.tex]{subfiles}

\begin{document}

\chapter{Variational Principle}

\label{chap:quantum_variational}
\noindent\rule{11cm}{0.4pt} \\
\textbf{Prerequisites:}  \autoref{chap:schrodinger} \\
\textbf{Difficulty Level:} ** \\
\noindent\rule{11cm}{0.4pt}
\vspace{5mm}

Suppose that $H$ is a Hamiltonian which is too complex to allow for direct solution of Schrodinger's equation. The variational principle provides a method to construct structured approximations to the solution of $H$ by leveraging a lower bound inequality on the energy of the system. 
%This method is very useful for 
\section{The Variational Principle}
The variational principle states that for any normalized wave function $\psi$, we have that

\begin{align}
E \leq \langle \psi | H | \psi \rangle 
\end{align}

Here $E$ is the ground state energy for the Hamiltonian. To see this, note that the eigenfunction of $H$ form a complete basis for the base Hilbert space $\mathcal{H}$. %\textcolor{red}{TODO: Why? What are the conditions under which this is true?}

Let us denote these wave functions as $\psi_i$ with associated energies $E_i$. Then we can expand any potential wave function as
\begin{align}
    \psi &= \sum_i c_i \psi_i
\end{align}
We can then compute
\begin{align}
    \langle \psi | H | \psi \rangle &= \Big \langle \sum_i c_i \psi_i \Big | H \Big | \sum_j c_j \psi_j \Big  \rangle \\
    &= \Big \langle \sum_i c_i \psi_i \Big | \sum_j c_j E_j \psi_j \Big \rangle \\
    &= \sum_{ij} c^*_i c_j E_j \langle \psi_i | \psi_j \rangle \\
    &= \sum_{i} |c_i|^2 E_i \\
    &\geq E_{\textrm{ground}}
\end{align}

By choosing different guesses for $\psi$ and computing $\langle \psi | H | \psi \rangle$ we can upper bound the ground state quantum energy of the system. In the case that our guess is not normalized, we replace the inequality above with the normalized version
\begin{align}
    E_{\textrm{ground}} \leq \frac{\langle \psi | H | \psi \rangle }{\langle \psi | \psi \rangle}
\end{align}

%\textcolor{red}{TODO: Complete. This should be fairly standard discussion of the variational principle with some references.}

\section{Variational Monte Carlo}

The basic variational approach can be applied analytically for simple systems, but larger quantum mechanical systems require the use of numerical techniques. Suppose that have a parameterized wave function

\begin{align}
    \psi(\theta)
\end{align}

where $\theta$ is some set of parameters. We write
\begin{align}
    E_\theta &= \frac{\langle \psi(\theta) | H | \psi(\theta) \rangle }{\langle \psi(\theta) | \psi(\theta) \rangle}
\end{align}
Note that the denominator can be written as
\begin{align}
   \langle \psi(\theta) | \psi(\theta) \rangle &= \int \psi^*(\theta) \psi(\theta) dx \\
   &= \int \|\psi(\theta)\|^2 dx
\end{align}
The numerator can be written as

\begin{align}
    \langle \psi(\theta) | H | \psi(\theta) \rangle &= \int \psi^*(\theta) H \psi(\theta) dx \\
    &= \int \|\psi(\theta)\|^2 \frac{H \psi(\theta)}{\psi(\theta)} dx
\end{align}
Combining, we can write
\begin{align}
    E_{\theta} &= \frac{\int \|\psi(\theta)\|^2 \frac{H\psi(\theta)}{\psi(\theta)} dx}{\int \|\psi(\theta)\|^2 dx} \\
    &= \int \frac{\|\psi(\theta)\|^2}{\int \|\psi(\theta)\|^2 dx} \frac{H\psi(\theta)}{\psi(\theta)} dx
\end{align}
We define
\begin{align}
    q(\theta) &= \frac{\|\psi(\theta)\|^2}{\int \|\psi(\theta)\|^2 dx}
\end{align}
and note that $q$ is a probability distribution. So we can write
\begin{align}
    E_\theta &= \int q(\theta) \frac{H \psi(\theta)}{\psi(\theta)} \\ 
    &= E_q \left [\frac{H \psi(\theta)}{\psi(\theta)} \right ]
\end{align}

By sampling from $q$, this expectation can be computed numerically
\begin{align}
    E_\theta \approx \frac{1}{N} \sum_i \frac{H\psi(\theta)(x_i)}{\psi(\theta)(x_i)}, \quad x_i \sim q(\theta)
\end{align}

By varying choice of $\theta$, we can improve the lower bound $E_\theta$. The choice of $\psi(\theta)$ is critical for these approximations. A simple choice is to do a direct factorization. For example, for a system with $N$ electrons,

\begin{align}
    \psi(X) &= \psi(x_1)\psi(x_2)\dotsc \psi(x_n)
\end{align}
A more sophisticated decomposition makes use of what is called a Jastrow factor
\begin{align}
    \psi(X) &= \exp \left ( \sum_{ij} u(r_{ij}) \right )
\end{align}

%\textcolor{red}{TODO: Draw a connection from here to variational autoencoders which draw on many of the same ideas.}

%\textcolor{red}{TODO: Add Physika code for variational monte carlo algorithm}
\end{document}